% ===== Mục 5 — Thảo luận (≈ 0,75 trang) =====
% Viết gì: mỗi ý 2–3 câu CÓ SỐ từ mục 4. Người: Hiếu + Thắng.
\subsection{Khi nào cổ điển đủ, khi nào cần học sâu}
\TODO{So dòng (i)/(ii) với (iii) ở Bảng~\ref{tab:ED2} và E-B1: cổ điển tìm vùng chữ được nhưng không phân lớp, nên phải dựa vào luật trên chữ OCR — lỗi OCR kéo theo lỗi phân lớp. Phần nào cổ điển vẫn đáng dùng (tiền xử lý khi ảnh xấu, theo E-A1)}

\subsection{Học sâu theo khối so với VLM đầu-cuối}
\TODO{Thắng: VLM mạnh ở đâu (không cần huấn luyện, đọc ngữ cảnh), yếu ở đâu (bịa tên/liều, chậm, phụ thuộc API, gửi ảnh y tế ra ngoài). Dẫn số từ Bảng~\ref{tab:ED2}}

\subsection{Hạn chế}
\begin{itemize}
  \item VAIPE là ảnh in máy, sạch, cùng một mẫu đơn; kết quả chưa được kiểm chứng trên ảnh chụp điện thoại thật hay mẫu đơn của bệnh viện khác. Bộ suy giảm tổng hợp (Mục~\ref{sec:items-gt}) chỉ mô phỏng một phần.
  \item Đáp án khoản thuốc được dẫn xuất bằng chính luật của \khoi{D}, chỉ soát bằng mắt \KQ{số ảnh} ảnh; E-D1 $=1{,}000$ vì thế chỉ chứng tỏ luật nhất quán, không chứng tỏ đáp án đúng tuyệt đối.
  \item Từ điển tên thuốc lấy từ tập huấn luyện; \KQ{\% OOV} tên thuốc kiểm tra nằm ngoài từ điển và không được hưởng lợi.
  \item \TODO{phần đã cắt vì thời gian: E-B3 (đường cong dữ liệu), E-C2 (tiền xử lý vùng cắt), E-D3; Tesseract / B cổ điển nếu không chạy}
\end{itemize}

\subsection{Hướng phát triển}
\TODO{1--2 câu: thu ảnh chụp thật để đánh giá; huấn luyện YOLO/VietOCR có tăng cường suy giảm; bỏ vùng \texttt{other} trước OCR để tăng tốc; từ điển thuốc quốc gia thay cho từ điển tập huấn luyện}
